\begin{table*}[!t]
    \centering
    \small
    \setlength{\tabcolsep}{4.5pt}
    \begin{tabular*}{\textwidth}{@{\extracolsep{\fill}}>{\raggedright\arraybackslash}m{0.45cm}>{\raggedright\arraybackslash}m{3.35cm}>{\raggedright\arraybackslash}m{1.70cm}>{\raggedright\arraybackslash}m{3.85cm}>{\raggedright\arraybackslash}m{3.25cm}>{\centering\arraybackslash}m{1.45cm}@{}}
        \toprule
        \textbf{ID} & \textbf{Risk} & \textbf{AI RMF} & \textbf{Evidence in this report} & \textbf{Control available} & \textbf{Residual} \\
        \midrule
        $T1$ & Adversarial masking: a periodic component injected on a candidate frequency is under-represented and escapes forecast-residual detection & \shortstack[l]{MAP\\MEASURE} & Precondition is $e^{\theta_{\mathrm{lock}}}>1$ (Model~B) with the site locations of $\kappa_S$ (Model~D2); the pathway itself needs the paired recovery ratio of Model~A. Not yet evaluated & Spectral monitor on the raw stream, upstream of tokenisation, at the enumerated candidate set & high \\
        \cmidrule(lr){1-6}
        $T2$ & Systematic blind spot with no adversary: a legitimate plant harmonic sits on a candidate frequency and is persistently under-represented & \shortstack[l]{MAP\\MEASURE\\MANAGE} & Candidate set is closed-form in $(f_s,P,S)$ and needs no fit; the cost at those sites is decided by $e^{\theta_{\mathrm{lock}}}$. Not yet evaluated & Choose $(P,S)$ against the known spectral content of the plant; \valueof{PartiallyObservable} degradation & medium \\
        \cmidrule(lr){1-6}
        $T3$ & Configuration drift: a change of sampling interval, patch size or checkpoint silently relocates the blind set & \shortstack[l]{GOVERN\\MAP} & $H3a$: relocation is predicted in closed form by $f_s/S$, and whether the sites follow it is decided by $\kappa_S$. Not yet evaluated & Recompute the candidate set from the knowledge base at every configuration change & low \\
        \cmidrule(lr){1-6}
        $T4$ & False assurance: overlap is increased and treated as a fix & \shortstack[l]{GOVERN\\MANAGE} & $M1$ is decided by $\delta_O$ under a two-leg rule and neither leg is evaluated. Independently of any fit, the patch branch is invariant to $S$ by construction & Do not credit overlap as a control; record the unestablished status in the model card & high \\
        \cmidrule(lr){1-6}
        $T5$ & Defence placed inside the blind spot: phase randomisation, or anomaly detection on forecast residuals & MANAGE & $H2$ is decided by $\sigma_\phi$ (Model~C): a deficit independent of phase makes randomising it useless. Not yet evaluated & Move detection upstream of tokenisation; treat residual-based detection as insufficient alone & medium \\
        \bottomrule
    \end{tabular*}
    \caption{Risk register. \textbf{AI RMF} names the function of the framework\autocite{tabassiArtificialIntelligenceRisk2023} under which each item is handled. \textbf{Residual} is the exposure that survives the control in the adjacent column, judged qualitatively: \textbf{high} where no control available to an operator closes the exposure, \textbf{medium} where a control exists but is partial or requires design-time choices, \textbf{low} where the exposure is closed by a procedure that follows from the arithmetic. Each entry names the estimand specified in \nameref{sec:six} that decides it; no posterior is asserted here, since the fits are not yet available. The two entries carried at \textbf{high} residual are those for which no control available to an operator closes the exposure: $T1$, because a detector placed after tokenisation inherits the blind spot it is meant to find, and $T4$, because the pre-registered mitigation cannot reach the patch branch by the arithmetic of \nameref{sec:eight} whatever the fits return.}
    \label{tab:riskRegister}
\end{table*}
